\documentclass[a4paper]{article}
\usepackage[a4paper,top=15mm,bottom=15mm,left=20mm,right=20mm]{geometry}
\usepackage{polyglossia}
\setdefaultlanguage[babelshorthands=true]{russian}
\setotherlanguage{english}
\usepackage{fontspec}
\setmainfont{FreeSerif}
\newfontfamily{\russianfonttt}[Scale=0.85]{DejaVuSansMono}
\usepackage[tiny,compact]{titlesec}
\usepackage{titling}
\setlength{\droptitle}{-1cm}
\pretitle{\begin{center}\begin{bfseries}\Large}
\posttitle{\par\end{bfseries}\end{center}}
\preauthor{\begin{center}\normalsize}
\postauthor{\par\end{center}\vspace{-1.5cm}}
\usepackage{hyperref}
\usepackage{bookmark}
\usepackage{booktabs}
\usepackage{float}

\title{Примитивы линейной алгебры: сравнительный анализ форматов хранения разреженных матриц}
\author{Семенский Роман Николаевич}
\date{}

\begin{document}
\maketitle
\begin{flushright}
Группа: \emph{25.Б42-мм} \\
Кафедра: \emph{системного программирования} \\
Научный руководитель: \emph{Григорьев Семён Вячеславович} \\
Номер семестра практики: \emph{2}
\end{flushright}

\section{Введение}
Работа выполнялась в рамках учебной практики и проекта по реализации примитивов разреженной линейной алгебры. Разреженная матрица содержит преимущественно нулевые элементы. В формате разреженного хранения записываются ненулевые элементы и их индексы.

В проекте создана собственная реализация формата COO (Coordinate List~--- список координат) на языке~C. Её производительность сравнивалась с CSparse~--- библиотекой для операций над разреженными матрицами, входящей в SuiteSparse.

Профилирование использовалось для определения наиболее затратных фаз алгоритмов и выбора направлений их оптимизации.

\section{Цель и задачи работы}
Цель работы~--- реализовать операции умножения для разреженных матриц в формате COO, исследовать их производительность в сравнении с CSparse и определить наиболее затратные действия в собственных алгоритмах.

Для достижения цели были поставлены следующие задачи:
\begin{enumerate}
\item реализовать умножение двух COO-матриц;
\item реализовать умножение COO-матрицы на разреженный вектор без его преобразования в плотный массив;
\item подготовить матрицы разного размера из SuiteSparse Matrix Collection;
\item измерить время двух операций для COO и CSparse по единой методике;
\item измерить затраты времени по фазам собственных алгоритмов.
\end{enumerate}

\section{Описание решения}
В формате COO для каждого ненулевого элемента хранятся индекс строки, индекс столбца и значение. Эти данные размещаются в трёх одномерных массивах. Для упорядочивания элементов по индексам строк и столбцов используется функция \texttt{qsort}.

При умножении матриц для второго сомножителя строится индекс строк. Значения одной строки результата накапливаются во временном массиве по индексам столбцов. Отдельная битовая карта отмечает столбцы, в которых выполнялось накопление. При формировании результата просматривается эта карта, а не весь временный массив.

Разреженный вектор представляется как COO-матрица размера $n\times1$. Перед умножением по его ненулевым элементам строится хеш-таблица. Она позволяет искать значение по индексу без двоичного поиска. Размер этой таблицы зависит от числа ненулевых элементов вектора, поэтому плотный вектор из $n$ значений не создаётся.

Для сравнения CSparse получает те же матрицы в формате CSC (Compressed Sparse Column~--- сжатое хранение по столбцам). Входные представления подготавливаются до замеров; измеряется только операция умножения.

\section{Эксперименты}
\subsection{Методика измерений}
В сравнительном эксперименте проверялась гипотеза о том, что на выбранном наборе матриц CSparse выполняет операции $A\times A$ и $A\times x$ быстрее собственной COO-реализации. Использовались девять матриц из SuiteSparse Matrix Collection. Набор включает матрицы размером от $62\times62$ до $9877\times9877$. Для симметричных файлов Matrix Market при загрузке явно добавляются элементы зеркальной части. Поэтому NNZ (Number of Non-Zeros~--- число ненулевых элементов) в таблицах указано после загрузки.

Измерялись $A\times A$ и $A\times x$. Вектор $x$ имел тот же размер, что и число столбцов $A$, а доля его ненулевых элементов составляла 10\%, поэтому вектор оставался разреженным. Перед созданием векторов генератор псевдослучайных чисел инициализировался фиксированным значением 42. Перед замерами выполнялись три прогревочных вызова. Затем собирались 15 выборок; число операций в каждой увеличивалось, пока её длительность не достигала 100~мс. В таблицах приведена медиана времени одного вызова. Разброс выборок оценивался интервалом от первого до третьего квартиля (Q1--Q3).

Характеристики вычислительной системы:
\begin{itemize}
\item операционная система: Ubuntu~26.04;
\item процессор: 13th~Gen Intel Core i5-13420H;
\item оперативная память: 16~ГБ;
\item компилятор: GCC (GNU Compiler Collection)~15.2.0, флаг оптимизации \texttt{-O3}.
\end{itemize}
Время измерялось монотонным таймером \texttt{clock\_gettime} с \texttt{CLOCK\_MONOTONIC}.

\subsection{Результаты сравнения}
В табл.~\ref{tab:mm} приведено время одного вызова $A\times A$. Последний столбец содержит отношение медианы времени COO к медиане времени CSparse.
\begin{table}[H]
\centering\small
\caption{Медиана времени умножения матриц $A\times A$}
\label{tab:mm}
\begin{tabular}{lrrrrr}
\toprule
Матрица & Размер & NNZ & COO, мкс & CSparse, мкс & COO/CSparse\\
\midrule
dolphins & $62\times62$ & 318 & 5{,}347 & 2{,}452 & 2{,}18\\
lesmis & $77\times77$ & 508 & 11{,}844 & 6{,}176 & 1{,}92\\
polbooks & $105\times105$ & 882 & 39{,}430 & 26{,}976 & 1{,}46\\
football & $115\times115$ & 1226 & 49{,}700 & 51{,}976 & 0{,}96\\
celegansneural & $297\times297$ & 2345 & 133{,}192 & 67{,}593 & 1{,}97\\
netscience & $1589\times1589$ & 5484 & 251{,}388 & 71{,}145 & 3{,}53\\
add20 & $2395\times2395$ & 17319 & 1535{,}335 & 1146{,}039 & 1{,}34\\
ca-GrQc & $5242\times5242$ & 28980 & 2749{,}474 & 1446{,}849 & 1{,}90\\
ca-HepTh & $9877\times9877$ & 51971 & 6724{,}213 & 3209{,}461 & 2{,}10\\
\bottomrule
\end{tabular}
\end{table}

Для $A\times A$ медиана времени CSparse была меньше на восьми матрицах. На матрице \texttt{football} медианы составили 49{,}700~мкс для COO и 51{,}976~мкс для CSparse. Интервалы Q1--Q3 для них равны 46{,}080--54{,}814~мкс и 49{,}791--63{,}258~мкс соответственно и пересекаются.

В табл.~\ref{tab:mv} приведены результаты $A\times x$. В обеих реализациях вектор оставался разреженным.
\begin{table}[H]
\centering\small
\caption{Медиана времени умножения матрицы на вектор $A\times x$}
\label{tab:mv}
\begin{tabular}{lrrrrr}
\toprule
Матрица & Размер & NNZ & COO, мкс & CSparse, мкс & COO/CSparse\\
\midrule
dolphins & $62\times62$ & 318 & 1{,}015 & 0{,}155 & 6{,}55\\
lesmis & $77\times77$ & 508 & 1{,}512 & 0{,}136 & 11{,}12\\
polbooks & $105\times105$ & 882 & 2{,}611 & 0{,}206 & 12{,}67\\
football & $115\times115$ & 1226 & 3{,}395 & 0{,}214 & 15{,}86\\
celegansneural & $297\times297$ & 2345 & 7{,}594 & 0{,}324 & 23{,}44\\
netscience & $1589\times1589$ & 5484 & 29{,}990 & 1{,}652 & 18{,}15\\
add20 & $2395\times2395$ & 17319 & 184{,}974 & 3{,}111 & 59{,}46\\
ca-GrQc & $5242\times5242$ & 28980 & 254{,}539 & 8{,}372 & 30{,}40\\
ca-HepTh & $9877\times9877$ & 51971 & 747{,}919 & 16{,}982 & 44{,}04\\
\bottomrule
\end{tabular}
\end{table}

Для $A\times x$ медиана времени CSparse была меньше на всех девяти матрицах. Отношение времени COO к времени CSparse находилось в пределах от 6{,}55 до 59{,}46.

Для сопоставления разброса использовалась относительная ширина интервала $(Q3-Q1)/M$, где $M$~--- медиана. В табл.~\ref{tab:spread} приведены наименьшие и наибольшие значения этой величины среди девяти матриц.
\begin{table}[H]
\centering\small
\caption{Относительная ширина интервала Q1--Q3}
\label{tab:spread}
\begin{tabular}{lrr}
\toprule
Операция и реализация & Минимум, \% & Максимум, \%\\
\midrule
$A\times A$, COO & 1{,}29 & 23{,}27\\
$A\times A$, CSparse & 1{,}75 & 25{,}91\\
$A\times x$, COO & 0{,}60 & 25{,}30\\
$A\times x$, CSparse & 1{,}29 & 26{,}17\\
\bottomrule
\end{tabular}
\end{table}
Значения Q1 и Q3 для каждой матрицы сохранены вместе с исходными результатами измерений; ссылка приведена в конце отчёта.

Для $A\times x$ проверяемая гипотеза подтвердилась на всех девяти матрицах. Для $A\times A$ медиана времени CSparse была меньше на восьми матрицах. На матрице \texttt{football} интервалы Q1--Q3 пересекались, поэтому имеющихся измерений недостаточно, чтобы утверждать о преимуществе одной из реализаций. Таким образом, на выбранном наборе матриц CSparse в большинстве случаев выполняла рассматриваемые операции быстрее собственной COO-реализации.

\subsection{Профилирование COO-реализации}
Цель профилирования~--- установить, какие действия в собственных алгоритмах вносят наибольший вклад во время выполнения. Измерялись проверка упорядоченности входных данных, построение индекса, накопление произведений, обработка битовой карты и создание результата. Профильная сборка использовала \texttt{-O3}.

Для профилирования $A\times A$ выбраны три самые крупные матрицы набора: выполнено 300 повторений на \texttt{add20}, 100 на \texttt{ca-GrQc} и 50 на \texttt{ca-HepTh}. Для $A\times x$ выбраны две самые крупные матрицы: выполнено 1000 повторений на \texttt{ca-GrQc} и 500 на \texttt{ca-HepTh}.

В исходной версии $A\times A$ после вычисления каждой строки просматривался весь временный массив. Эта фаза занимала 76{,}66\% времени на \texttt{add20}, 86{,}22\% на \texttt{ca-GrQc} и 90{,}23\% на \texttt{ca-HepTh}. В исходной версии $A\times x$ двоичный поиск и накопление занимали 93{,}31\% на \texttt{ca-GrQc} и 94{,}63\% на \texttt{ca-HepTh}. Следовательно, дальнейшие изменения были направлены на исключение полного просмотра временного массива и повторного двоичного поиска.

После профилирования полный просмотр временного массива был заменён обработкой битовой карты затронутых столбцов, а двоичный поиск~--- поиском в хеш-таблице. Цель повторного эксперимента~--- проверить, сокращают ли эти изменения время выполнения собственных алгоритмов. В табл.~\ref{tab:optimization} приведены медианы времени одного вызова до и после изменений. Они получены по методике сравнительного эксперимента, описанной выше. Снижение рассчитано относительно исходной версии.
\begin{table}[H]
\centering\small
\caption{Результат оптимизации COO-реализации}
\label{tab:optimization}
\begin{tabular}{llrrr}
\toprule
Операция & Матрица & До, мкс & После, мкс & Снижение, \%\\
\midrule
$A\times A$ & add20 & 2882{,}296 & 1535{,}335 & 46{,}73\\
$A\times A$ & ca-GrQc & 13094{,}578 & 2749{,}474 & 79{,}00\\
$A\times A$ & ca-HepTh & 43480{,}173 & 6724{,}213 & 84{,}53\\
$A\times x$ & ca-GrQc & 1175{,}114 & 254{,}539 & 78{,}34\\
$A\times x$ & ca-HepTh & 2499{,}207 & 747{,}919 & 70{,}07\\
\bottomrule
\end{tabular}
\end{table}

После внесённых изменений медиана времени уменьшилась на всех матрицах, выбранных для повторного эксперимента. Для $A\times A$ снижение составило 46{,}73--84{,}53\%, для $A\times x$~--- 70{,}07--78{,}34\%. Следовательно, предположение о существенном влиянии полного просмотра временного массива и повторного двоичного поиска на время выполнения подтвердилось. Исходные данные и программа профилирования указаны в конце отчёта.

\section{Заключение}
В ходе работы были получены следующие результаты.
\begin{itemize}
\item Реализовано умножение двух COO-матриц.
\item Реализовано умножение COO-матрицы на разреженный вектор без преобразования в плотный массив.
\item Для экспериментов подготовлены девять матриц размером от $62\times62$ до $9877\times9877$ из SuiteSparse Matrix Collection.
\item Сравнительный эксперимент показал, что на выбранном наборе матриц CSparse имела меньшую медиану времени в 17 из 18 сравнений. Для $A\times A$ на матрице \texttt{football} интервалы Q1--Q3 пересекались, поэтому преимущество одной из реализаций в этом случае не установлено.
\item Профилирование показало, что значительная часть времени выполнения приходилась на полный просмотр временного массива при $A\times A$ и повторный двоичный поиск при $A\times x$. Их замена обработкой отмеченных элементов и поиском в хеш-таблице уменьшила медиану времени на всех матрицах, выбранных для повторного эксперимента.
\end{itemize}

\section*{Ссылки}
\begin{itemize}
\item исходный код и скрипты измерений: \url{https://github.com/rmnsmnsk/implementation-of-linear-algebra-primitives} (дата обращения: 12.09.2026);
\item техническая документация: \url{https://github.com/rmnsmnsk/implementation-of-linear-algebra-primitives/blob/main/README.md} (дата обращения: 12.09.2026);
\item результаты сравнительных измерений: \url{https://github.com/rmnsmnsk/implementation-of-linear-algebra-primitives/blob/main/benchmark_results_wsl.txt} (дата обращения: 12.09.2026);
\item исходные данные профилирования: \url{https://github.com/rmnsmnsk/implementation-of-linear-algebra-primitives/blob/main/profiling_results.csv} (дата обращения: 12.09.2026);
\item сравнение времени до и после оптимизации: \url{https://github.com/rmnsmnsk/implementation-of-linear-algebra-primitives/blob/main/optimization_results.csv} (дата обращения: 12.09.2026);
\item исходные результаты профилирования: \url{https://github.com/rmnsmnsk/implementation-of-linear-algebra-primitives/blob/main/profiling_raw_wsl.txt} (дата обращения: 12.09.2026);
\item программа профилирования: \url{https://github.com/rmnsmnsk/implementation-of-linear-algebra-primitives/blob/main/src/profile.c} (дата обращения: 12.09.2026);
\item SuiteSparse Matrix Collection: \url{https://sparse.tamu.edu/} (дата обращения: 12.09.2026);
\item репозиторий SuiteSparse и CSparse: \url{https://github.com/DrTimothyAldenDavis/SuiteSparse} (дата обращения: 12.09.2026);
\item репозиторий с отчётами: \url{https://github.com/gsvgit/1_year_practice_2026_reports} (дата обращения: 12.09.2026).
\end{itemize}
\end{document}
